\section{Exact algebraic contracts}
\subsection{Repeated-root-safe domain support}
Let $D\subseteq K$ be finite and let
$L_D(X)=\prod_{d\in D}(X-d)$. The empty product is one.

\begin{theorem}[Support by nilpotence]\label{thm:support}
For a monic $H\in K[X]$ of degree $m\geq1$ and an integer $k\geq m$,
every root of $H$ in an algebraic closure lies in $D$ if and only if
\[
          (L_D\bmod H)^k=0\quad\text{in }K[X]/(H).
\]
\end{theorem}
\begin{proof}
Write $H=\prod_\alpha(X-\alpha)^{e_\alpha}$ over an algebraic closure.
The displayed equality is equivalent to $H\mid L_D^k$. The roots of $L_D$
are precisely the elements of $D$, each with multiplicity one. Divisibility
therefore holds exactly when every $\alpha\in D$ and every $e_\alpha\leq k$.
The second condition follows from $e_\alpha\leq m\leq k$.
Divisibility over the algebraic closure is equivalent to divisibility over
$K$ by the remainder theorem for monic division.
\end{proof}

The nonconstant hypothesis matters: a constant monic $H=1$ has no roots and
must be handled as a separate trivial case. If $D$ is empty, the criterion
rejects every nonconstant $H$. Squarefreeness of $H$ is not assumed.
In characteristic $p$ one may take $k=p^{\lceil\log_p m\rceil}$.
For a binary quintic, $k=8$ suffices, as in the compact support oracle \ev{E1}.
This supplies an exact predicate, not a bound on how many coefficient
candidates must be searched.

For a binary linear space $V$ and $v\notin V$,
\[
 L_{V+\F_2v}(X)=L_V(X)L_V(X+v)
              =L_V(X)^2+L_V(v)L_V(X).
\]
Indeed the two cosets are disjoint, and induction starting at $L_{\{0\}}=X$
shows that $L_V$ is additive. For an affine space $a+V$, substitute $X-a$.
These formulas can represent an exponentially large binary domain implicitly.
They do not give a proper nonzero $\F_p$-linear subspace of the prime field
$\F_p$; an interval or another domain representation is required there.

In an extension field, Frobenius acts on coefficients too:
$(\sum c_iX^i)^p=\sum c_i^pX^{ip}$. It is $\F_p$-linear, generally not
$K$-linear. For example, in $\F_4$ with $\omega^2=\omega+1$,
the square of the constant $\omega$ is $\omega+1$, not $\omega$.
In odd characteristic an ordinary square also contains cross terms.
Neither binary shortcut can be transplanted unchanged to an odd-field backend.

\subsection{A binary elliptic-curve norm chart}
Throughout this subsection $K=\F_{2^n}$ and
\[
 E:\quad Y^2+XY=C(X),\qquad C(X)=X^3+aX^2+b,\qquad b\ne0.
\]
Fix $T=(z,y_T)\in E(K)$ with $z\ne0$. For $r,s,h,k\in K$ set
\begin{equation}
 \begin{split}
 P&=X^3+rX^2+sX+t,\qquad Q=hX+k,\\
 t&=z^3+rz^2+sz+(hz+k)y_T,\\
 N&=P^2+XPQ+Q^2C .
 \end{split}
\end{equation}
Conjugation in the quadratic function-field equation sends $Y$ to $Y+X$,
so $N$ is the norm of $P+QY$. The chosen $t$ makes $(P+QY)(T)=0$;
hence $N(z)=0$. Its leading coefficient is one and its degree is six.
Thus $H=N/(X+z)$ is a monic quintic.

Given ordered source values $x_1,\ldots,x_5\in K$, write
\[
 H_x=\prod_{i=1}^5(X+x_i)
     =X^5+e_1X^4+e_2X^3+e_3X^2+e_4X+e_5.
\]
Let $\beta=\sqrt b$, $c=(y_T+\beta)/z$, and $w=\sqrt{e_5/z}$.
Square roots are unique because Frobenius is an automorphism of $K$.

\begin{theorem}[Exact norm-chart parametrization]\label{thm:chart}
The identity $N=(X+z)H_x$ is equivalent to the following conditions:
\begin{align}
 h^2+h&=e_1+z,                                             \label{eq:as}\\
 k&=e_2+r^2+hr+ah^2+ze_1,                                 \label{eq:k}\\
 s&=zr+ck+y_Th+z^2+w,                                     \label{eq:s}\\
 0&=R_3:=e_3+hs+kr+k^2+ze_2,                              \label{eq:r3}\\
 0&=R_4:=\beta k^2+(zw)k+e_5+ze_4.                        \label{eq:r4}
\end{align}
No hypothesis excludes $h=0$, $k=0$, zero sources, or repeated sources.
\end{theorem}
\begin{proof}
The coefficients of $N$, in descending order, are
\[
\begin{array}{c|l}
6&1\\
5&h+h^2\\
4&r^2+k+hr+ah^2\\
3&kr+hs+k^2\\
2&s^2+ks+ht+bh^2+ak^2\\
1&kt\\
0&t^2+bk^2 .
\end{array}
\]
Comparison at degrees five, four, and three gives
\eqref{eq:as}, \eqref{eq:k}, and \eqref{eq:r3}.
The constant coefficient identity is $t^2+bk^2=ze_5$, or
$t=\beta k+zw$; substituting the definition of $t$ and dividing by the
nonzero $z$ gives \eqref{eq:s}. The coefficient at degree one becomes
\eqref{eq:r4}. This proves necessity.

Conversely these equations match every coefficient except possibly degree two.
The difference $N+(X+z)H_x$ is therefore $AX^2$ for some $A\in K$.
Both polynomials vanish at $z$, giving $Az^2=0$. Since $z\ne0$, $A=0$.
This proves sufficiency without another division or a generic-position
assumption.
\end{proof}

The map $\AS(h)=h^2+h$ is $\F_2$-linear with kernel $\{0,1\}$ and
image dimension $n-1$. Choose any fixed linear section $S$ on this image.
For arbitrary inputs $e_1$, extend the section linearly to $K$ and check
the explicit residual \eqref{eq:as}. Then
\[
                  h=S(e_1+z)+\eta,\qquad \eta\in\F_2
\]
enumerates each valid $h$ exactly once. This works for even as well as odd
$n$; an odd-degree half-trace formula is not an all-extension implementation.
Once the sources, $r$, and $\eta$ are fixed, \eqref{eq:k} and
\eqref{eq:s} determine $k,s$ uniquely. The explicit residuals then accept
exactly the norm identities. For a $d$-dimensional binary source space,
the primary description has $5d+n+1$ bits \ev{E1--E3}; the actual
$n=83,d=18$ fixture uses 174.

This bijection concerns norm witnesses with the specified target and degree
chart. The target branch $z=0$ is outside the theorem. A norm factor can
correspond to the conjugate point above a source coordinate, and special
divisors and chain auxiliary fibers require separate analysis. The theorem
does not imply that these 174-bit models are in bijection with the original
339-variable chain solutions.
